% Фрагмент графа способов получения объектов: путь к железной кирке.
\begin{tikzpicture}[
  font=\footnotesize,
  thing/.style={draw, align=center, minimum width=2.5cm, minimum height=8mm, inner sep=2pt},
  part/.style={-{Stealth[length=1.8mm]}, semithick},
  tool/.style={-{Stealth[length=1.8mm]}, semithick, dashed},
  station/.style={-{Stealth[length=1.8mm]}, semithick, dotted},
  lbl/.style={font=\scriptsize, inner sep=1.5pt, fill=white},
]

\node[thing] (log) at (0, 0) {Бревно};
\node[thing] (planks) at (3.3, 0) {Доски};
\node[thing] (stick) at (6.6, 0) {Палка};
\node[thing] (woodpick) at (0, -1.7) {Деревянная\\кирка};
\node[thing] (cobble) at (3.3, -1.7) {Булыжник};
\node[thing] (stonepick) at (6.6, -1.7) {Каменная\\кирка};
\node[thing] (furnace) at (9.9, -1.7) {Печь};
\node[thing] (ore) at (6.6, -3.4) {Железная\\руда};
\node[thing] (ingot) at (9.9, -3.4) {Железный\\слиток};
\node[thing] (ironpick) at (13.2, -1.7) {Железная\\кирка};

\draw[part] (log) -- node[lbl] {1} (planks);
\draw[part] (planks) -- node[lbl] {2} (stick);
\draw[part] (planks.south west) -- node[lbl] {3} (woodpick.north east);
\draw[part] (cobble) -- node[lbl] {3} (stonepick);
\draw[part] (stick) -- node[lbl] {2} (stonepick);
\draw[part] (ore) -- node[lbl] {1} (ingot);
\draw[part] (ingot.east) -| node[lbl, pos=0.75] {3} (ironpick.south);
\draw[part] (stick.east) -| node[lbl, pos=0.75] {2} (ironpick.north);
\draw[tool] (woodpick) -- (cobble);
\draw[tool] (stonepick) -- (ore);
\draw[station] (furnace) -- (ingot);

% --- условные обозначения ---
\draw[part] (0.2, -4.6) -- ++(1.1, 0) node[right, font=\scriptsize] {входит в состав, число~--~количество};
\draw[tool] (8.4, -4.6) -- ++(1.1, 0) node[right, font=\scriptsize] {требуется для добычи};
\draw[station] (0.2, -5.15) -- ++(1.1, 0) node[right, font=\scriptsize] {оборудование рецепта};

\end{tikzpicture}
